%%%% Начало оформления заголовка - оставить без изменений !!! %%%%
\input{my_folder/task_settings}	% настройки - начало 
	
				{%\normalfont %2020
						\MakeUppercase{\SPbPU}}\\
				\institute

\par}\intervalS% завершает input

				\noindent
				\begin{minipage}{\linewidth}
				\vspace{\mfloatsep} % интервал 	
				\begin{tabularx}{\linewidth}{Xl}
					&УТВЕРЖДАЮ      \\
					&Руководитель образовательной про-\\
					&граммы «Прикладная математика и \\
					&информатика»     \\			
					&\underline{\hspace*{0.1\textheight}} \Head     \\
					&<<\underline{\hspace*{0.05\textheight}}>> \underline{\hspace*{0.1\textheight}} \approveYear г.  \\  
				\end{tabularx}
				\vspace{\mfloatsep} % интервал 	
				\end{minipage}

\intervalS{\centering\bfseries%

				ЗАДАНИЕ\\
				на выполнение %с 2020 года 
				%по выполнению % до 2020 года
				выпускной квалификационной работы


\intervalS\normalfont%

				студенту \uline{Губину Р.В.{} гр.~\group}


\par}\intervalS%
%%%%
%%%% Конец оформления заголовка  %%%%
 	
	
	
\begin{enumerate}[1.]
	\item Тема работы: {\expandafter \ulined \thesisTitle}
	%\item Тема работы (на английском языке): \uline{\thesisTitleEn.} % вероятно после 2021 года
	\item Срок сдачи студентом законченной работы: \uline{\thesisDeadline.} 
	\item Исходные данные по работе: \\
	Система передачи данных с 8 и более абонентами\\
	Скорость передачи данных между абонентами -- от 32 бит/с до 37 кбит/с\\
	Инструментальные средства: \\
	•\hspace{1cm} Язык программирования Python \\
	Ключевые источники литературы : 
	\begin{enumerate}[1.]
	\item Компьютерные сети. Нисходящий подход.  Куроуз Д.Ф., Росс К.В. 2016 
	\item Stallings, William. Data and Computer Communications. 10th ed., Pearson, 2015. 
	\item "Python Network Programming Cookbook", R. Kumar, 2017 
	\item "Fluent Python: Сlеаг, Concise Programming", Luciano Ramalho, 2022.
	\end{enumerate}
	\item Содержание работы (перечень подлежащих разработке вопросов):
	\begin{enumerate}[label=\arabic*)]
		\item Введение. Обоснование актуальности проблемы 
		\item Постановка задачи. Формулировка функции цели и ограничений. 
		\item Обзор методов моделирования низкоскоростных систем передачи данных.\\
		Уравнения и численные методы.
		\item Разработка модели. Настройка параметров модели по измерениям и\\
		верификация.
		\item Проведение численных экспериментов и анализ результатов. 
		\item Выводы
		\item Заключение
	\end{enumerate}
		\item Дата выдачи задания: \uline{\thesisStartDate.}
\end{enumerate}

\intervalS%можно удалить пробел

Руководитель ВКР \uline{\hrulefill}\Supervisor

\intervalS%можно удалить пробел

%Консультант по нормоконтролю \uline{\hspace*{0.1\textheight} \ConsultantNorm}%ПОКА НЕ ТРЕБУЕТСЯ, Т.К. ОН У ВСЕХ ПО УМОЛЧАНИЮ

Задание принял к исполнению 

\intervalS%можно удалить пробел

Студент \uline{\hrulefill}\Author



\input{my_folder/task_settings_restore}	% настройки - конец